\chapter{Syntax Quick Reference}
\label{app:quickref}

A condensed index of the constructs used in this guide, grouped the way you
meet them: declarations first, then iteration, tile selection and movement,
then statements, coordination and diagnostics. Chapter~\ref{ch:programming-guide}
explains what each construct means and why it is defined that way; this
appendix is for looking up a spelling without re-reading an example.

\section{Declarations and shapes}

\begin{table}[htbp]
\centering
\begin{tabularx}{\linewidth}{lX}
\toprule
Construct & Meaning \\
\midrule
\co@__co__ T [shape] f(T [shape] a)@ & Croqtile function with shaped parameters. \\
\co@__co__ void@ / \co@__co__ auto@ & In-place result / inferred return type. \\
\co@T [a.span] b;@ & Buffer with a copied shape. \\
\co@a.span@ / \co@a.span(i)@ & Whole shape / one dimension. \\
\co@x : [7, 8]@ & Named \co@mdspan@ declaration. \\
\co@a = {1, 2, 3}@ & \co@ituple@ declaration. \\
\co@tensor.at(i, j)@ & Position access. \\
\bottomrule
\end{tabularx}
\end{table}

\section{Iteration}

\begin{table}[htbp]
\centering
\begin{tabularx}{\linewidth}{lX}
\toprule
Construct & Meaning \\
\midrule
\co@parallel p by N@ / \co@parallel {p, q} by [M, N]@ & Independent instances. \\
\co@: block@ / \co@: thread@ / \co@: group@ & Annotations on a parallel axis. \\
\co@foreach idx in [N]@ / \co@foreach t(1:)@ & Ordered / sliced iteration. \\
\co@with idx in [a, b]@ & Scoped index binding. \\
\co@p # m@ / \co@#p@ & Compose indices / axis extent. \\
\co@cdiv(a, b)@ & Ceiling division. \\
\bottomrule
\end{tabularx}
\end{table}

\section{Tile selection and movement}

\begin{table}[htbp]
\centering
\begin{tabularx}{\linewidth}{lX}
\toprule
Construct & Meaning \\
\midrule
\co@f = dma.copy src => shared@ & Tile movement statement. \\
\co@a.chunkat(t)@ & One tile of a regular partition. \\
\co@a.subspan(M, K).at(i, j)@ & Explicit extents and tile coordinates. \\
\co@a.subspan(M, K).step(sM, sK).at(i, j)@ & Strided tile selection. \\
\co@a.view(M, N).from(r, c)@ & Window with explicit element offsets. \\
\co@a.span_as([dims])@ & Shape reinterpretation without moving data. \\
\co@f.data@ / \co@f.span@ / \co@dma.any@ & Moved tile / its shape / placeholder handle. \\
\co@.async@ / \co@wait f@ / \co@swap@ / \co@rotate@ & Async transfer and handle rotation. \\
\co@global@ / \co@shared@ / \co@local@ & Storage specifiers. \\
\bottomrule
\end{tabularx}
\end{table}

\section{Statement modifiers and contraction}

\begin{table}[htbp]
\centering
\begin{tabularx}{\linewidth}{lX}
\toprule
Construct & Meaning \\
\midrule
\co@dma.transp<...>@ / \co@dma.pad<...>@ & Transform while moving. \\
\co@.zfill@ / \co@.swiz<N>@ & Fill and layout modifiers. \\
\co@mc = mma.fill 0.0@ & Zero an accumulator. \\
\co@ma = mma.load src.chunkat(...)@ & Load an operand. \\
\co@mma.row.row mc, ma, mb@ & Contraction with a layout contract. \\
\co@mma.store mc, dst@ & Store the accumulator. \\
\co@.sp@ / \co@.scale@ / \co@mma.commit@ / \co@mma.scale@ & Contraction variants and stage fence. \\
\bottomrule
\end{tabularx}
\end{table}

\section{Coordination, interop and diagnostics}

\begin{table}[htbp]
\centering
\begin{tabularx}{\linewidth}{lX}
\toprule
Construct & Meaning \\
\midrule
\co@if (expr) { ... }@ & Runtime predicate. \\
\co@inthreads.async (cond)@ / \co@inthreads (cond)@ & Concurrent region / sequential filter. \\
\co@shared event e[N]@ / \co@trigger e@ / \co@wait e@ & Synchronisation tokens. \\
\co@parallel.async ...@ / \co@stream s@ / \co@yield@ & Host scheduling and early return. \\
\co@__device__ fn()@ / \co@call fn(args)@ & Interoperating C++ function and its invocation. \\
\co@__cpp__("...")@ / \co@__cpp__(R"(...)")@ & Injected source text. \\
\co@#define@ / \co@#if@ / \co@#ifdef@ / \co@#error@ / \co@-D@ & Preprocessor. \\
\co@print!@ / \co@println!@ / \co@print@ / \co@println@ & Compile-time and runtime output. \\
\co@assert@ / \co@croq_assert@ & Device and host assertions. \\
\bottomrule
\end{tabularx}
\end{table}

\begin{brief}[Where to go next]
For the meaning behind a construct, return to the section that introduced it:
Chapter~\ref{ch:language-basics} for the first reading and
Chapter~\ref{ch:programming-guide} for the design. To see a construct in a
complete program, start with the worked kernels in
Section~\ref{sec:3-4-1}, \ref{sec:3-6-1} and \ref{sec:3-8-5}.
\end{brief}
